\section{Summary of Contributions}
This thesis set out to determine whether trimodal cross-modal attention fusion could detect and characterise the structure of affective incoherence---the systematic divergence between voluntary emotional expression (audio, video) and involuntary neural response (EEG)---at population scale. Both the fusion performance objective and the coherence-characterisation objective were achieved on the EAV corpus. The \texttt{TrimodalAttentionFusion} module reached 80.2\% mean test accuracy across 42 subjects, and softhard modality-dropout training of the same architecture reached 84.7\% mean test accuracy under full modality and 81.5\% under the zero-EEG audio-visual configuration that corresponds to a realistic deployment scenario. More substantively, the $5\times5$ suppression matrix produced four interpretable, cross-population patterns of affective incoherence---bidirectional Anger--Happiness confusion and directional Sadness/Calmness$\rightarrow$Neutral suppression---validated on a sample of 5{,}040 test trials without the EEG-Neutral inflation artifact that had appeared in the 3-subject pilot.

This thesis set out to determine whether trimodal cross-modal attention fusion could detect and characterise the structure of affective incoherence---the systematic divergence between voluntary emotional expression (audio, video) and involuntary neural response (EEG)---at population scale, and to establish what such a system achieves on participants it has never seen. Under the within-subject protocol the \texttt{TrimodalAttentionFusion} module reached 80.2\% mean test accuracy across 42 subjects, rising to 84.7\% with softhard modality-dropout training and 81.5\% under the zero-EEG audio-visual configuration. Under five-fold subject-wise cross-validation the same pipeline reaches 62.2\%, and the ordering of fusion methods reverses. More substantively, the $5\times5$ suppression matrix produced four interpretable, cross-population patterns of affective incoherence---bidirectional Anger--Happiness confusion and directional Sadness/Calmness$\rightarrow$Neutral suppression---validated on 5{,}040 test trials.

\textbf{Contribution 1: Trimodal fusion architecture combining cross-modal attention, modality-dropout robustness, and coherence analysis on EAV.} Three trimodal fusion methods have previously been published on the EAV benchmark: AMERL by Yin et al.\ \cite{Yin2024AMERL} at 70.86\%, Hyper-MML by Kang et al.\ \cite{Kang2025HyperMML} at 76.65\% (the prior published state of the art), and EEG-MoCE \cite{Anon2026EEGMoCE} at 75.88\%. None of these methods includes a missing-modality robustness mechanism, a zero-EEG inference path, or a cross-modal coherence analysis. The present work is, to our knowledge, the first to integrate all three of these capabilities on the EAV benchmark within a single architecture. On the per-subject within-subject evaluation protocol used by the open-source comparator (AMERL), the present work reaches 80.2\% (cross-attention fusion) and 84.7\% (with softhard modality dropout), establishing a new state-of-the-art trimodal fusion result on EAV under that protocol, and additionally publishes the five reference numbers---naive late fusion (77.5\%), cross-attention fusion (80.2\%), concat-MLP fusion (81.7\%), modality-dropout fusion (84.7\%), and zero-EEG demo path (81.5\%)---against which future EAV fusion work can be compared. All numbers are reproducible from the open-source codebase at \texttt{shafinmuffinn/thesis\_p2}.

\textbf{Contribution 2: Systematic identification and resolution of EAV repository bugs.} Five bugs in the EAV open-source repository were identified and resolved: a forward-hook return-value defect in EEGNet that caused systematic training failure, a train/eval mode cycling defect in the shared trainer, a classifier-head mismatch in the vision trainer, a RAM-exhaustion failure in vision preprocessing, and a dependency conflict introduced by \texttt{facenet-pytorch}. The EEGNet bugs in particular---which caused EEG accuracy to be chance-level in the repository's published state---have implications for any researcher attempting to reproduce or extend the EAV paper's EEGNet baseline using the PyTorch path.

\textbf{Contribution 3: Trimodal cross-modal attention fusion module.} The \texttt{TrimodalAttentionFusion} module is an open-source, tested, and extensible trimodal fusion architecture specifically designed for the EAV modality dimensions. It satisfies two architectural commitments: it returns per-modality embeddings as named dictionary outputs, enabling future temporal head attachment; and it produces finite predictions when EEG is zeroed, enabling inference in EEG-unavailable deployment scenarios.

\textbf{Contribution 4: First suppression matrix characterisation on EAV.} The $5\times5$ suppression matrix computed over 5,040 test trials constitutes the first population-scale characterisation of cross-modal affective coherence on EAV. Three dominant patterns are identified---bidirectional Anger--Happiness high-arousal confusion, directional Sadness$\rightarrow$Neutral suppression, and directional Calmness$\rightarrow$Neutral suppression---all confirmed as cross-population phenomena. This matrix provides a reference baseline against which future coherence-detection methods on EAV can be compared.

\textbf{Contribution 5: Reframing of the EAV evaluation task.} The thesis's central methodological contribution is the reframing of multimodal emotion recognition on EAV from a pure accuracy-maximisation problem to a coherence-characterisation problem. This reframing is motivated by the epistemological asymmetry between EEG (involuntary, internal) and audio-video (voluntary, external), and by the clinical significance of cross-modal divergence. It provides a research agenda for future EAV-based work that goes beyond asking which architecture achieves the highest accuracy and instead asks what cross-modal disagreement reveals about affective regulation.

\textbf{Contribution 6: First subject-independent evaluation of trimodal fusion on EAV.} All prior trimodal-fusion work on this benchmark---AMERL \cite{Yin2024AMERL}, Hyper-MML \cite{Kang2025HyperMML}, and EEG-MoCE \cite{Anon2026EEGMoCE}---reports within-subject results only. Chapter~\ref{chap:crosssubject} reports a five-fold subject-wise cross-validation over all 42 participants, yielding three findings without precedent on this corpus. First, the headline accuracy falls from 84.7\% to 62.2\%, quantifying what the conventional protocol conceals. Second, the loss is localised: audio transfers at parity ($+0.4$ pp) while vision collapses by 33.4 percentage points, and training the vision encoder longer makes its cross-subject accuracy worse---evidence that EAV's strong within-subject vision results are substantially identity-driven. Third, the architectural ordering reverses: naive late fusion significantly outperforms every learned fusion variant ($p < 0.003$), and the learned variants become statistically indistinguishable from one another. These findings bear on the interpretation of every within-subject result published on this benchmark, including those of this thesis.

\section{Limitations}
Four limitations are explicitly acknowledged.

\textbf{Limitation 1: Controlled corpus, not spontaneous data.} The EAV corpus uses cue-based emotion elicitation. The ``suppression events'' detected in the matrix may partly reflect imperfect emotion simulation, where a participant correctly displays the cued emotion in audio-video but fails to enter the corresponding physiological state fully, rather than genuine social suppression. This confound is irreducible without cross-referencing against a concurrent internal-state measure.

\textbf{Limitation 2: Corpus-bounded generalisation.} The within-subject limitation of the original evaluation is retired by Chapter~\ref{chap:crosssubject}, which reports subject-independent results over all 42 participants. A narrower limitation remains: all results, under both protocols, are obtained on a single laboratory corpus recorded under one protocol with one hardware configuration. Generalisation across corpora, recording conditions, and populations is untested, and the identity-inflation finding for vision is established on EAV alone.

\textbf{Limitation 3: Single pooled-vector feature resolution.} The per-modality encoders produce single pooled vectors rather than token sequences. This limits the cross-modal information accessible to the attention fusion module and likely suppresses the attention module's potential advantage over naive averaging. Replacing the pooled-vector encoders with sequence-producing encoders, such as applying LSTM or TCN modules to sequences of per-frame ViT representations, is the primary architectural improvement available. The cross-subject results give this limitation additional weight: under subject shift the three learned fusion variants are statistically indistinguishable from one another, which is what would be expected if the feature resolution, rather than the fusion operator, were the binding constraint.

\textbf{Limitation 4: No clinical validation.} The suppression matrix findings are validated behaviorally through cross-population consistency, stable base rate, and absence of the EEG-Neutral inflation confound, but not clinically. The patterns may or may not correspond to clinically meaningful suppression in healthy participants, and the methodology has not been applied to any clinical population. Any clinical inference from the present results would be premature.

\section{Future Work}
Three workstreams are identified as Phase 2 extensions, each directly building on the architectural and methodological infrastructure developed in this thesis.

\textbf{Phase 2.1: Clinical population validation.} The suppression-matrix methodology, applied to a population with a known affective dysregulation diagnosis such as alexithymia measured by the Toronto Alexithymia Scale, would test whether the cross-modal coherence score discriminates diagnostic groups from healthy controls. This is the primary scientific value of the methodology: the EAV-based characterisation provides a reference baseline, and deviations from that baseline in clinical populations are the signal of interest.

\textbf{Phase 2.2: Primary audio-visual dataset and merged training.} A primary recording study would collect audio and video from 5--10 new subjects performing the same five-emotion task as EAV, recorded with consumer-grade smartphones in non-laboratory settings. The primary dataset would be merged with EAV through two-stage fine-tuning: the AST and ViT encoders are pre-trained on EAV's audio and speaking-clip vision data, then fine-tuned on the union of EAV speaking clips and the new primary recordings with appropriate per-corpus weighting. The EEG branch would continue to train on EAV alone, with the primary dataset contributing only audio-visual examples to the fusion module. At inference on primary recordings, the EEG token is zeroed, matching the demo-path configuration of Section~\ref{sec:modality_dropout_results}. This design tests two properties not testable on EAV alone: cross-corpus generalisation of the audio-visual encoders, and the ecological validity of the zero-EEG demo path on data recorded outside the EAV laboratory protocol. Cross-modal coherence on the primary dataset would substitute self-reported internal-state Likert ratings, collected per trial, for the unavailable EEG signal; cross-modal disagreement is then measured between audio-vision consensus and self-report. EEG instrumentation, ethics approval for primary EEG recording, and behavioural-EEG validation of the Likert proxy are required for a full trimodal extension that includes a primary EEG channel.

\textbf{Phase 2.3: In-the-wild temporal emotion tracking.} The architecture designed in this thesis---with per-modality embeddings exposed as named outputs---supports attachment of a temporal head, such as an LSTM or Residual-TCN, for continuous emotion tracking on temporally annotated video corpora such as DFEW or MAFW. This extension requires a temporally annotated in-the-wild corpus, a voice activity detector for audio, a robust face detector and tracker for video, and a temporal decoder consuming the per-modality embedding sequences. The per-modality encoders (AST, ViT, EEGNet) can be frozen, making this an incremental addition rather than a full redesign.

\textbf{Phase 2.4: MERCL contrastive pre-training.} The supervised contrastive pre-training scheme of Lee, Kim and Kim \cite{Lee2024MERCL}---comprising intra-modal contrastive loss (AMCL), inter-modal alignment loss (EMCL), and sample-wise alignment (SMCL)---was dropped from the present thesis under schedule pressure. Its re-introduction in Phase 2 would test whether cross-modal contrastive alignment of the per-modality encoder representations, before fine-tuning the fusion head, materially improves either classification accuracy or the signal-to-noise ratio of the suppression matrix. The hypothesis, based on the Lee et al. ablation results, is that contrastive pre-training meaningfully improves cross-modal feature alignment and would manifest as a larger fusion gain on EAV. Concretely, AMCL would cluster intra-modal class representations, EMCL would align same-class representations across modalities, and SMCL would close the modality-specific embedding gap with a margin term $\alpha$. All three losses operate on the pre-classifier embeddings already exposed by the \texttt{forward()} contract, requiring no architectural changes.

\textbf{Phase 2.5: Subject-invariant visual representation.} The identity inflation documented in Section~\ref{sec:cs_permodality} identifies a concrete target. If a substantial portion of the vision encoder's within-subject accuracy is attributable to participant identity, then explicitly discouraging identity information---through adversarial identity confusion, identity-invariant contrastive objectives, or per-participant feature normalisation---should improve cross-subject accuracy even at some cost to within-subject accuracy. This is directly testable with the protocol and infrastructure already in place: the measurement apparatus exists, and the cross-subject baseline of 41.2\% is the number to beat.

\textbf{Closing remark.} The thesis opens and closes with the same asymmetry: EEG is involuntary; audio and video are not. This asymmetry, usually a confound in fusion system design, is here the signal. The suppression matrix transforms that signal into a population-level map of affective regulation. Whether that map reflects genuine emotion suppression, simulation artifacts, or cross-modal latency mismatches remains an open question---one that the behavioural evidence in this thesis positions cleanly for future clinical investigation.